\documentclass{labreport}

\usepackage{microtype}
\usepackage{enumitem}
\usepackage{pifont}
\usepackage{tabularx}
\usepackage{booktabs}

% Headheight adjustment for fancyhdr
\setlength{\headheight}{28pt}
\addtolength{\topmargin}{-13pt}

% Colors
\definecolor{codebg}{RGB}{248, 249, 251}
\definecolor{codeborder}{RGB}{218, 222, 229}
\definecolor{linkcolor}{RGB}{20, 60, 120}
\definecolor{zshgreen}{RGB}{0, 160, 80}
\definecolor{zshred}{RGB}{200, 30, 30}

% Hyperref setup with subtle link colors
\hypersetup{
    colorlinks=true,
    linkcolor=linkcolor,
    filecolor=linkcolor,
    urlcolor=linkcolor,
    citecolor=linkcolor
}

% Code and terminal listings
\lstdefinestyle{bash}{
    basicstyle=\ttfamily\footnotesize,
    backgroundcolor=\color{codebg},
    frame=single,
    framerule=0.5pt,
    rulecolor=\color{codeborder},
    breaklines=true,
    breakatwhitespace=false,
    showstringspaces=false,
    keepspaces=true,
    columns=flexible,
    numbers=none,
    xleftmargin=3pt,
    xrightmargin=3pt,
    framesep=4pt
}

\lstdefinestyle{c}{
    language=C,
    basicstyle=\ttfamily\footnotesize,
    backgroundcolor=\color{codebg},
    frame=single,
    framerule=0.5pt,
    rulecolor=\color{codeborder},
    keywordstyle=\color{blue!85!black}\bfseries,
    commentstyle=\color{gray!80!black}\itshape,
    stringstyle=\color{teal!80!black},
    numbers=left,
    numberstyle=\tiny\color{gray!80!black},
    numbersep=8pt,
    breaklines=true,
    breakatwhitespace=false,
    showstringspaces=false,
    keepspaces=true,
    columns=flexible,
    tabsize=2,
    xleftmargin=14pt,
    xrightmargin=3pt,
    framesep=4pt
}

\lstset{
    extendedchars=true,
    literate=
        {➜}{{\textcolor{zshgreen}{\ding{220}}}}{1}
        {├─}{|-\,}{2}
        {└─}{\`-\,}{2}
        {├}{|}{1}
        {─}{-}{1}
        {└}{\`}{1}
        {✗}{{\textcolor{zshred}{\ding{55}}}}{1}
}

% Section styling for clean document structure
\titleformat{\section}
  {\large\bfseries}{\thesection}{0.8em}{}[\vspace{1pt}\hrule height 0.5pt]
\titleformat{\subsection}
  {\normalsize\bfseries}{\thesubsection}{0.8em}{}
\titleformat{\subsubsection}
  {\small\bfseries}{\thesubsubsection}{0.8em}{}

\begin{document}

\labtitle
    {9}
    {Watching System Calls with strace}
    {st. gr. FAF-243}
    {Șoimu Ionuț}
    {asist. univ.\\[0.15cm] Reițman Patricia}

\tableofcontents
\newpage

% ==============================================================================
\section{Lab 9: Track A -- Watching System Calls with strace (Linux)}

\subsection{Before you start}

\subsubsection*{Purpose and Scope}
In operating systems theory, the boundary between an unprivileged user-space program and the privileged kernel represents the most critical architectural partition of modern computing. User-space programs execute in restricted hardware modes (Ring 3 on x86\_64) where direct interaction with hardware devices, physical memory page tables, and scheduling queues is strictly prohibited. To access persistent storage, create network sockets, or communicate with the display, a program must request services through the \textbf{system call interface}.

The objective of Laboratory Work 9 is to inspect, intercept, and analyze this user/kernel threshold from user space using the diagnostic utility \texttt{strace}. By monitoring both standard utility binaries (such as GNU \texttt{ls}) and custom low-level C programs invoking raw POSIX system calls (\texttt{open}, \texttt{write}, \texttt{read}, \texttt{close}), we observe:
\begin{enumerate}[label=(\alph*), noitemsep]
    \item How user-space library calls translate into kernel trap instructions;
    \item How dynamic linkers, runtime libraries, and environment lookups contribute to the syscall footprint of seemingly trivial commands;
    \item How the kernel allocates file descriptor integers according to POSIX standards;
    \item How kernel diagnostic tracers capture execution frequencies, latency distributions, and system call error states.
\end{enumerate}

\subsubsection*{Tooling and Environment}
Experiments were carried out on a 64-bit Linux installation running Linux Kernel 6.x. The primary investigative tool is \texttt{strace}, which utilizes the kernel's \texttt{ptrace(2)} system call mechanism to halt a monitored child process at every syscall entry and exit point, inspecting CPU registers and memory buffers to decode call names, parameters, and return codes.

% ==============================================================================
\subsection{Part 1. See the calls of a real command}

\subsubsection{Monitoring Command Execution with strace}
Standard command-line utilities frequently execute extensive initialization routines before reaching their core application logic. To observe this in practice, we executed GNU \texttt{ls} under \texttt{strace}, filtering specifically for basic I/O and file descriptor system calls:

\begin{lstlisting}[style=bash]
➜  os-lab9 strace -f -e trace=openat,read,write,close ls -l
\end{lstlisting}

\subsubsection{Observed Terminal Session}
Below is the trace output captured during the execution of \texttt{ls -l}:

\begin{lstlisting}[style=bash]
openat(AT_FDCWD, "/etc/ld.so.cache", O_RDONLY|O_CLOEXEC) = 3
close(3)                                = 0
openat(AT_FDCWD, "/usr/lib/libcap.so.2", O_RDONLY|O_CLOEXEC) = 3
read(3, "\177ELF\2\1\1\0\0\0\0\0\0\0\0\0\3\0>\0\1\0\0\0\200\212\0\0\0\0\0\0"..., 1024) = 1024
close(3)                                = 0
openat(AT_FDCWD, "/usr/lib/libc.so.6", O_RDONLY|O_CLOEXEC) = 3
read(3, "\177ELF\2\1\1\3\0\0\0\0\0\0\0\0\3\0>\0\1\0\0\0P\177\2\0\0\0\0\0"..., 1024) = 1024
close(3)                                = 0
openat(AT_FDCWD, "/usr/lib/locale/locale-archive", O_RDONLY|O_CLOEXEC) = 3
close(3)                                = 0
openat(AT_FDCWD, "/usr/share/locale/locale.alias", O_RDONLY|O_CLOEXEC) = 3
read(3, "# Locale name alias data base.\n#"..., 4096) = 2998
read(3, "", 4096)                       = 0
close(3)                                = 0
openat(AT_FDCWD, "/etc/nsswitch.conf", O_RDONLY|O_CLOEXEC) = 3
read(3, "# Name Service Switch configurat"..., 4096) = 359
read(3, "", 4096)                       = 0
close(3)                                = 0
openat(AT_FDCWD, "/etc/passwd", O_RDONLY|O_CLOEXEC) = 3
read(3, "root:x:0:0::/root:/usr/bin/bash\n"..., 4096) = 2245
close(3)                                = 0
openat(AT_FDCWD, "/etc/group", O_RDONLY|O_CLOEXEC) = 3
read(3, "root:x:0:root\nsys:x:3:bin\nmem:x:"..., 4096) = 1218
close(3)                                = 0
openat(AT_FDCWD, "/home/ionut/Programming/OS/os-lab9", O_RDONLY|O_NONBLOCK|O_CLOEXEC|O_DIRECTORY) = 3
close(3)                                = 0
write(1, "total 68\n", 9)               = 9
openat(AT_FDCWD, "/etc/localtime", O_RDONLY|O_CLOEXEC) = 3
read(3, "TZif2\0\0\0\0\0\0\0\0\0\0\0\0\0\0\0\0\0\0\20\0\0\0\20\0\0\0\0"..., 4096) = 2424
read(3, "TZif2\0\0\0\0\0\0\0\0\0\0\0\0\0\0\0\0\0\0\20\0\0\0\20\0\0\0\0"..., 4096) = 1509
close(3)                                = 0
write(1, "-rw-r--r-- 1 ionut ionut 65991 Nov  9 13:48 LOGO_UTM.jpg\n", 57) = 57
write(1, "drwxr-xr-x 1 ionut ionut   226 Nov  9 13:48 report\n", 48) = 48
close(1)                                = 0
close(2)                                = 0
+++ exited with 0 +++
\end{lstlisting}

\subsubsection{Analysis of Calls}
Even for a simple listing command, the operating system executes a substantial sequence of system calls:
\begin{itemize}[noitemsep]
    \item \textbf{Dynamic Linker Resolution:} The process opens `/etc/ld.so.cache` to locate shared libraries (`libcap.so.2`, `libc.so.6`), verifies their ELF headers via `read()`, and maps them into its virtual memory.
    \item \textbf{Identity and Name Service Mapping:} In order to display textual user and group names (`ionut ionut`), \texttt{ls} consults `/etc/nsswitch.conf`, `/etc/passwd`, and `/etc/group`.
    \item \textbf{Locale and Timezone:} Timestamps require timezone rules extracted from `/etc/localtime`.
    \item \textbf{Directory Traversal and Output:} The targeted directory descriptor is opened via `openat()`, directory entries are retrieved, and formatted rows are written line-by-line to standard output.
\end{itemize}
In total, a single execution of `ls -l` makes between \textbf{50 and 70 system calls}.

\subsubsection*{Observe:}
\begin{enumerate}[leftmargin=*]
    \item \textbf{Which system call does ls use to print to the screen? To which descriptor does it write?}\\[0.2em]
    \texttt{ls} uses the \textbf{\texttt{write()}} system call to print to the screen, and it writes specifically to \textbf{file descriptor \texttt{1}} (\texttt{STDOUT\_FILENO}, standard output).\\
    As seen in the trace output:
    \begin{lstlisting}[style=bash]
write(1, "total 68\n", 9) = 9
write(1, "-rw-r--r-- 1 ionut ionut 65991 Nov  9 13:48 LOGO_UTM.jpg\n", 57) = 57
    \end{lstlisting}
    Each line of the directory listing is handed directly to the kernel through descriptor 1.
\end{enumerate}

% ==============================================================================
\subsection{Part 2. A program that calls the kernel directly}

\subsubsection{Direct Kernel Invocation: syscalls.c}
In high-level C programming, developers typically invoke standard I/O library functions such as \texttt{fopen()}, \texttt{fprintf()}, and \texttt{fgets()}. These library functions maintain internal user-space buffers in an attempt to minimize expensive context switches into supervisor mode.

To observe direct kernel transitions without runtime buffering, we implemented \texttt{syscalls.c}, which interacts with persistent storage solely through raw POSIX system calls:

\begin{lstlisting}[style=c]
/* syscalls.c - open, write, read and close, using only system calls.
 * Build: gcc -Wall -Wextra -o syscalls syscalls.c
 * Run:   ./syscalls
 * See the calls: strace ./syscalls
 */
#include <fcntl.h>      /* open */
#include <unistd.h>     /* write, read, close */
#include <string.h>     /* strlen */
#include <stdio.h>      /* printf */

int main(void) {
    const char *msg = "hello, kernel\n";

    /* open() is a system call: ask the kernel for a file, get back a descriptor */
    int fd = open("note.txt", O_CREAT | O_WRONLY | O_TRUNC, 0644);
    if (fd < 0) { perror("open"); return 1; }
    printf("open returned descriptor %d\n", fd);

    /* write() is a system call: hand the bytes to the kernel */
    ssize_t n = write(fd, msg, strlen(msg));
    printf("write sent %zd bytes\n", n);
    close(fd);

    /* read it back, through the kernel again */
    fd = open("note.txt", O_RDONLY);
    char buf[64];
    n = read(fd, buf, sizeof(buf) - 1);
    buf[n] = '\0';
    printf("read got %zd bytes: %s", n, buf);
    close(fd);
    return 0;
}
\end{lstlisting}

\subsubsection{Compilation and Direct Execution}
We compiled and executed the program without tracing:

\begin{lstlisting}[style=bash]
➜  os-lab9 gcc -Wall -Wextra -o syscalls syscalls.c
➜  os-lab9 ./syscalls
open returned descriptor 3
write sent 14 bytes
read got 14 bytes: hello, kernel
\end{lstlisting}

\subsubsection{Tracing System Calls with strace}
Next, we executed the compiled binary under \texttt{strace}, filtering specifically for \texttt{openat}, \texttt{write}, \texttt{read}, and \texttt{close}:

\begin{lstlisting}[style=bash]
➜  os-lab9 strace -e trace=openat,write,read,close ./syscalls
openat(AT_FDCWD, "/etc/ld.so.cache", O_RDONLY|O_CLOEXEC) = 3
close(3)                                = 0
openat(AT_FDCWD, "/usr/lib/libc.so.6", O_RDONLY|O_CLOEXEC) = 3
read(3, "\177ELF\2\1\1\3\0\0\0\0\0\0\0\0\3\0>\0\1\0\0\0P\177\2\0\0\0\0\0"..., 1024) = 1024
close(3)                                = 0
openat(AT_FDCWD, "note.txt", O_WRONLY|O_CREAT|O_TRUNC, 0644) = 3
write(1, "open returned descriptor 3\n", 27) = 27
write(3, "hello, kernel\n", 14)         = 14
write(1, "write sent 14 bytes\n", 20)   = 20
close(3)                                = 0
openat(AT_FDCWD, "note.txt", O_RDONLY)  = 3
read(3, "hello, kernel\n", 63)          = 14
write(1, "read got 14 bytes: hello, kernel\n", 33) = 33
close(3)                                = 0
+++ exited with 0 +++
\end{lstlisting}

\subsubsection{Correlation Between Code and Kernel Traps}
The trace reveals a strict 1-to-1 correspondence between our C source statements and the emitted kernel traps:
\begin{enumerate}
    \item \texttt{open("note.txt", O\_CREAT | O\_WRONLY | O\_TRUNC, 0644)} triggers:
    \begin{center}
        \texttt{openat(AT\_FDCWD, "note.txt", O\_WRONLY|O\_CREAT|O\_TRUNC, 0644) = 3}
    \end{center}
    On modern 64-bit Linux kernels, libc implements \texttt{open()} by calling the generalized \texttt{openat()} system call passing \texttt{AT\_FDCWD} (representing the current working directory).
    \item \texttt{printf("open returned descriptor \%d\textbackslash n", fd)} flushes to stdout:
    \begin{center}
        \texttt{write(1, "open returned descriptor 3\textbackslash n", 27) = 27}
    \end{center}
    \item \texttt{write(fd, msg, strlen(msg))} writes 14 bytes to the file:
    \begin{center}
        \texttt{write(3, "hello, kernel\textbackslash n", 14) = 14}
    \end{center}
    \item \texttt{close(fd)} releases descriptor 3:
    \begin{center}
        \texttt{close(3) = 0}
    \end{center}
    \item \texttt{open("note.txt", O\_RDONLY)} reopens the file:
    \begin{center}
        \texttt{openat(AT\_FDCWD, "note.txt", O\_RDONLY) = 3}
    \end{center}
    \item \texttt{read(fd, buf, sizeof(buf) - 1)} reads up to 63 bytes, receiving 14:
    \begin{center}
        \texttt{read(3, "hello, kernel\textbackslash n", 63) = 14}
    \end{center}
    \item \texttt{close(fd)} terminates the file session:
    \begin{center}
        \texttt{close(3) = 0}
    \end{center}
\end{enumerate}

\subsubsection*{Observe:}
\begin{enumerate}[leftmargin=*]
    \item \textbf{The descriptor returned by the first openat is 3, not 0, 1 or 2. Why are those three already taken?}\\[0.2em]
    Under the POSIX system standard, every standard process inherits three pre-allocated file descriptors at process instantiation:
    \begin{itemize}[noitemsep, topsep=2pt]
        \item \textbf{Descriptor 0 (\texttt{STDIN\_FILENO}):} Standard Input stream.
        \item \textbf{Descriptor 1 (\texttt{STDOUT\_FILENO}):} Standard Output stream.
        \item \textbf{Descriptor 2 (\texttt{STDERR\_FILENO}):} Standard Error diagnostic stream.
    \end{itemize}
    The Linux kernel's file descriptor allocator adheres to the deterministic \textbf{lowest-available-descriptor rule}: whenever a new descriptor is requested via \texttt{open()}, \texttt{creat()}, \texttt{dup()}, or \texttt{pipe()}, the kernel scans the process's file descriptor table starting from index 0 and allocates the smallest integer slot that is not currently in use.\\
    Since slots 0, 1, and 2 are already occupied by the terminal streams established by the shell, the first user file descriptor allocated is guaranteed to be integer \textbf{3}.
\end{enumerate}

% ==============================================================================
\subsection{Part 3. Count one kind of call}

\subsubsection{Statistical Profiling with strace -c}
The \texttt{-c} flag enables profiling mode in \texttt{strace}. Instead of printing each system call sequentially, \texttt{strace} attaches timers to each call and presents a statistical summary upon process exit.

We profiled the execution of \texttt{./syscalls} filtering specifically for \texttt{write}:

\begin{lstlisting}[style=bash]
➜  os-lab9 strace -c -e write ./syscalls
open returned descriptor 3
write sent 14 bytes
read got 14 bytes: hello, kernel
% time     seconds  usecs/call     calls    errors syscall
------ ----------- ----------- --------- --------- ----------------
100.00    0.000039           9         4           write
------ ----------- ----------- --------- --------- ----------------
100.00    0.000039           9         4           total
\end{lstlisting}

\subsubsection{Deconstructing the 4 write Calls}
Although our source code contains only one explicit call to \texttt{write()} (\texttt{write(fd, msg, strlen(msg))}), the profiler reports exactly \textbf{4 write system calls}.
The discrepancy is explained by the three calls to \texttt{printf()}:
\begin{enumerate}[noitemsep]
    \item \texttt{write(1, "open returned descriptor 3\textbackslash n", 27)} -- produced by \texttt{printf}
    \item \texttt{write(3, "hello, kernel\textbackslash n", 14)} -- our explicit system call
    \item \texttt{write(1, "write sent 14 bytes\textbackslash n", 20)} -- produced by \texttt{printf}
    \item \texttt{write(1, "read got 14 bytes: hello, kernel\textbackslash n", 33)} -- produced by \texttt{printf}
\end{enumerate}
Because standard output connected to a terminal is line-buffered by default in \texttt{glibc}, every newline character (\texttt{\textbackslash n}) in a \texttt{printf()} format string triggers a flush, generating an underlying \texttt{write(1, ...)} system call.

\subsubsection{Complete System Call Profiling for ./syscalls}
Running an unfiltered profiling session (\texttt{strace -c ./syscalls}) demonstrates the total syscall load of a minimal C program:

\begin{lstlisting}[style=bash]
➜  os-lab9 strace -c ./syscalls
open returned descriptor 3
write sent 14 bytes
read got 14 bytes: hello, kernel
% time     seconds  usecs/call     calls    errors syscall
------ ----------- ----------- --------- --------- ----------------
 28.99    0.000060          15         4           close
 22.71    0.000047          11         4           openat
 15.94    0.000033           8         4           write
 10.63    0.000022           7         3           mprotect
  8.70    0.000018           9         2           getrandom
  5.80    0.000012           4         3           brk
  3.38    0.000007           3         2           read
  1.93    0.000004           1         3           fstat
  1.93    0.000004           4         1           prlimit64
  0.00    0.000000           0         8           mmap
  0.00    0.000000           0         1         1 access
  0.00    0.000000           0         1           execve
  0.00    0.000000           0         1           arch_prctl
  0.00    0.000000           0         1           set_tid_address
  0.00    0.000000           0         2           newfstatat
  0.00    0.000000           0         1           set_robust_list
  0.00    0.000000           0         1           rseq
------ ----------- ----------- --------- --------- ----------------
100.00    0.000207           4        42         1 total
\end{lstlisting}

The program executed \textbf{42 system calls} in 207 microseconds:
\begin{itemize}[noitemsep]
    \item Memory management calls (\texttt{mmap}, \texttt{mprotect}, \texttt{brk}) allocate stack guards, map the C library, and initialize the heap.
    \item Process control calls (\texttt{execve}, \texttt{arch\_prctl}, \texttt{set\_tid\_address}, \texttt{rseq}) configure CPU thread-local storage and scheduling parameters.
    \item File I/O calls (\texttt{openat}, \texttt{read}, \texttt{write}, \texttt{close}) handle library loading, user data transfer, and stream cleanup.
\end{itemize}

% ==============================================================================
\section{Conclusion}

Through Laboratory Work 9, I investigated the system call interface of the Linux operating system from the outside using \texttt{strace}.

Analyzing the execution of standard GNU utilities revealed that even elementary commands such as \texttt{ls -l} require dozens of system calls to resolve shared dynamic libraries, translate numerical user/group IDs to names via `/etc/passwd` and `/etc/group`, decode locale configuration, and traverse directory trees. We confirmed that screen output is achieved exclusively via the \texttt{write()} system call directed to standard output file descriptor 1.

By implementing \texttt{syscalls.c}, I observed the direct correlation between user-level POSIX primitives and kernel trap dispatches. We verified that libc translates \texttt{open()} to modern \texttt{openat()} calls, and proved that descriptor allocation strictly follows the lowest-available-descriptor rule, causing the first user-opened file to receive descriptor 3 because descriptors 0, 1, and 2 are pre-allocated for standard input, standard output, and standard error.

Finally, profiling system calls with \texttt{strace -c} demonstrated how user-space runtime buffering in libc bridges high-level functions like \texttt{printf()} with kernel calls like \texttt{write()}. This laboratory provided concrete empirical confirmation of how the operating system kernel protects and multiplexes hardware resources behind uniform, traceable system call boundaries.

\end{document}
